%% sections/Research_Background.tex
\section{Related Work}\label{sec:research_background}

\subsection{\textbf{2. Literature Background and
Conservation-Oriented Problem
Definition}}\label{literature-background-and-conservation-oriented-problem-definition}

Existing research has made important progress in the digital
documentation and conservation of Dong drum towers. Previous studies
have examined their architectural typologies, roof-crown forms,
construction techniques, structural components, master-carpenter
knowledge, acoustic characteristics, mathematical modeling, UAV-based
recording, BIM reconstruction, and HBIM-oriented intelligent modeling.
These studies have provided valuable foundations for understanding Dong
drum towers as complex timber heritage systems and have demonstrated the
feasibility of using digital technologies to record, visualize, and
interpret their architectural characteristics. However, much of the
existing work has focused on documentation workflows, geometric
modeling, digital representation, or heritage information management.
Less attention has been paid to how different neural reconstruction
methods perform under comparable empirical conditions and how their
outputs can be interpreted in relation to conservation-oriented
documentation needs.

Neural reconstruction methods are relevant to Dong drum tower
documentation because these structures present challenges that are
difficult to address through conventional image-based modeling alone.
Their multi-layered eaves, dense timber frameworks, repeated roof
modules, roof-crown details, and shadowed interior spaces require
methods that can handle complex appearance, partial occlusion, viewpoint
variation, and fine architectural details. NeRF-based methods provide a
mature baseline for view synthesis through continuous radiance-field
representation. INGP improves practicality by reducing training cost
through multiresolution hash encoding. 3DGS offers efficient explicit
scene representation and real-time rendering potential, making it
suitable for visually complex heritage scenes. 2DGS, although not
primarily optimized for image-level metrics in the present study,
remains relevant because of its surface-oriented formulation and
potential value for future geometry-focused evaluation. These
complementary characteristics justify a controlled comparison of NeRF,
2DGS, 3DGS, and INGP for Dong drum tower reconstruction and provide the
basis for the following discussion of conservation-relevant components
and related neural reconstruction methods.

\subsection{2.1 Conservation-Relevant Architectural Components of Dong
Drum
Towers}\label{conservation-relevant-architectural-components-of-dong-drum-towers}

For Dong drum towers, conservation-oriented documentation should focus
on architectural components that carry structural, morphological, and
cultural significance. The most critical elements include layered eaves,
roof-crown systems, columns and beams, bracket-like supporting members,
mortise-and-tenon joint areas, enclosure components, base structures,
and shadowed inner timber spaces. These components should not be
understood merely as visual features. Together, they encode the
construction logic, craft knowledge, spatial hierarchy, morphological
identity, and maintenance history of the tower.

Layered eaves are among the most visually distinctive and
conservation-relevant features of Dong drum towers. They reflect the
vertical hierarchy of the structure and the carpentry logic used to
organize repeated roof layers. However, their dense arrangement can
cause self-occlusion, repeated-texture ambiguity, and shadow
accumulation during image-based reconstruction. Roof-crown systems are
also important because they function as symbolic and morphological
markers of the tower, yet they are often small, elevated, and difficult
to capture with sufficient detail from ground-based or limited UAV
viewpoints. Timber columns, interlocking beams, bracket-like supports,
and mortise-and-tenon joint areas are essential for understanding the
structural configuration and traditional construction techniques, but
they are frequently affected by low illumination, dense wooden
intersections, and limited camera visibility.

The lower base, enclosure elements, inner timber spaces, and areas
beneath multi-layered eaves are also important for restoration survey
and periodic monitoring because they may contain traces of deformation,
material aging, moisture effects, or later repair interventions. These
areas are usually more difficult to reconstruct than exposed exterior
surfaces because of shadow, narrow viewing angles, and partial
occlusion. Therefore, a conservation-oriented evaluation of neural
reconstruction methods should not rely solely on global visual metrics.
It should also consider whether the reconstructed outputs preserve the
readability of critical architectural components and whether the results
can support practical tasks such as visual documentation, HBIM
preparation, restoration planning, digital exhibition, and long-term
monitoring.

These conservation-relevant components define the main technical
challenges for digital reconstruction. Classical registration and
model-fitting methods, such as ICP and RANSAC, provide essential support
for 3D alignment and geometric robustness, but they do not directly
address the view-synthesis and rendering challenges posed by complex
heritage scenes [28,29]. Existing-building BIM reconstruction and
large-scale 3D modeling studies have shown that complex heritage
environments often require the integration of point-cloud processing,
semantic enrichment, and object-based interpretation [30--32].
Knowledge-based HBIM enrichment and mesh-based information modeling
further indicate that high-quality reconstruction is not only a visual
task, but also a prerequisite for semantic representation and downstream
conservation use [33,34]. For Dong drum towers, these challenges are
amplified by multi-layered eaves, repeated roof geometries, dense timber
components, roof-crown details, and heavily shaded spaces beneath the
eaves.

\subsection{2.2 Related Work on Neural Reconstruction and Heritage
Documentation}\label{related-work-on-neural-reconstruction-and-heritage-documentation}

Recent advances in neural rendering have opened new possibilities for
the digital reconstruction of complex heritage scenes. Local Light Field
Fusion provided an early practical framework for view synthesis with
prescribed sampling strategies [35]. Neural Radiance Fields
demonstrated that scene geometry and appearance can be jointly encoded
as a continuous radiance field for high-quality novel-view synthesis
[36]. Mip-NeRF and Mip-NeRF 360 further improved anti-aliasing and
unbounded-scene rendering, making neural scene representation more
robust in real-world environments [37,38]. Instant-NGP greatly
reduced training cost through multiresolution hash encoding, improving
the practicality of neural rendering workflows [39]. More recently,
3D Gaussian Splatting achieved real-time radiance-field rendering with
strong visual fidelity [40], while 2D Gaussian Splatting emphasized
a more geometrically oriented surface representation [41].
Extensions such as GS-SLAM and deformable 3D Gaussians have further
demonstrated the broader applicability of Gaussian-based representations
to dense SLAM and dynamic-scene reconstruction [42,43].

In image-based rendering tasks, PSNR, SSIM, and LPIPS are widely used to
measure pixel-level fidelity, structural similarity, and perceptual
similarity, respectively [44,45]. These metrics provide a useful
basis for comparing rendered views with reference images, but they do
not fully capture the conservation value of reconstructed outputs. For
architectural heritage, especially complex timber structures,
reconstruction quality must also be interpreted in relation to component
readability, structural continuity, and downstream documentation tasks.
Research on 3D point clouds has similarly highlighted the importance of
geometry-aware interpretation and downstream spatial understanding
[46].

Despite these advances, the applicability of neural reconstruction
methods to complex timber heritage remains insufficiently examined. Most
existing studies on Dong drum towers have focused on sound-field
simulation, mathematical modeling, UAV-based documentation, or
HBIM-oriented modeling, rather than controlled cross-method evaluation
under the same empirical heritage dataset [9--12]. As a result, it
remains unclear how NeRF-, Gaussian-, and hash-encoding-based methods
behave when processing multi-eave timber structures with repeated
components, occluded spaces, and strong illumination variation. More
importantly, it remains unclear how quantitative rendering metrics,
computational efficiency, and architectural-component readability should
be jointly interpreted for practical method selection in digital
heritage workflows.

To move beyond visual demonstration toward conservation-oriented
evaluation, this study addresses three heritage-science-oriented
research questions. First, which neural reconstruction method most
reliably preserves conservation-relevant architectural components of
Dong drum towers, including layered eaves, roof-crown systems, timber
frameworks, and shadowed internal spaces? Second, how do different
neural scene representations---NeRF-based, Gaussian-based, and
hash-encoding-based---perform across heterogeneous drum tower scenes
acquired under comparable UAV imaging conditions? Third, how can
rendering fidelity, computational efficiency, and architectural-view
readability be integrated to guide task-oriented method selection for
digital documentation, rapid field reconstruction, HBIM preparation,
digital exhibition, and periodic monitoring?

To answer these questions, this study constructs a multi-case empirical
corpus using UAV images of four representative Dong drum towers: Chaoli,
Congjiang, Zeli, and Zengchong. Four neural rendering and reconstruction
methods---NeRF, 2DGS, 3DGS, and INGP---are implemented and compared
under a controlled experimental workflow based on consistent scene data,
SfM-based camera pose estimation, image screening, and train/test
splitting where applicable. Rendering quality is evaluated using PSNR,
SSIM, and LPIPS, while computational practicality is assessed through
training time and rendering speed. In addition, representative
architectural views, including the base, enclosure, main body, tower
tip, and panoramic view, are analyzed to connect metric-based evaluation
with conservation-oriented interpretation.

The main contributions of this study are fourfold. First, it establishes
a UAV-based multi-case empirical corpus for representative Dong drum
towers, providing a comparative basis for neural reconstruction research
in complex timber heritage contexts. Second, it presents a controlled
evaluation of NeRF, 2DGS, 3DGS, and INGP using consistent scene data,
camera poses, and quantitative rendering metrics. Third, it combines
PSNR, SSIM, LPIPS, training efficiency, rendering speed, and
architectural-view analysis to assess method performance from both
computational and heritage-documentation perspectives. Fourth, it
proposes a conservation-oriented interpretation framework that
translates algorithmic results into task-aware method-selection guidance
for digital visualization, rapid reconstruction, HBIM-oriented modeling,
restoration survey, and long-term monitoring of timber heritage.

Overall, this study shifts the digital reconstruction of Dong drum
towers from isolated visual demonstration toward a task-aware evaluation
framework for heritage conservation. By linking neural rendering
performance with the documentation needs of complex timber heritage, the
study provides practical evidence for selecting appropriate
reconstruction methods in conservation-oriented workflows and
contributes to the broader application of neural reconstruction
technologies in heritage science.

